% Generated draft: requires mathematical and source review.
\documentclass[11pt]{article}
\usepackage[margin=25mm]{geometry}
\usepackage{fontspec}
\setmainfont{Latin Modern Roman}
\usepackage{amsmath,amssymb,mathtools,unicode-math}
\setmathfont{Latin Modern Math}
\usepackage{enumitem,needspace,xurl,hyperref}
\setlength{\parindent}{0pt}
\setlength{\parskip}{5pt}
\setlength{\emergencystretch}{3em}
\newcommand{\N}{\mathbb{N}}
\newcommand{\Z}{\mathbb{Z}}
\newcommand{\Q}{\mathbb{Q}}
\newcommand{\R}{\mathbb{R}}
\newcommand{\C}{\mathbb{C}}
\newcommand{\F}{\mathbb{F}}
\newcommand{\A}{\mathbb{A}}
\newcommand{\PP}{\mathbb{P}}
\newcommand{\E}{\mathbb{E}}
\newcommand{\eps}{\varepsilon}
\newcommand{\dd}{\,\mathrm{d}}
\begin{document}
\section*{UnsolvedMath ID: 30004016}
\textbf{Source problem code:} catalogue entry
\subsection{Definitions and mathematical statement}
The supplied question is: \emph{How far can Jantzen's sum formula determine the characters of simple representations, beyond the groups where it is already complete?}

For conditional notation, suppose that an algebraic group \(G\) over a field \(k\), a split torus \(T\subseteq G\), and a finite-dimensional rational \(G\)-module \(V\) have been specified. Write \(X^*(T)=\operatorname{Hom}(T,\mathbb G_m)\) for the character lattice and \(V_\mu\) for the \(\mu\)-weight space of \(V\). Its formal weight character is
\[
\operatorname{ch}_T(V)=\sum_{\mu\in X^*(T)}\dim_k(V_\mu)e^\mu\in\mathbb Z[X^*(T)],
\]
where \(e^\mu\) denotes the basis element indexed by \(\mu\) in the group ring. These definitions describe one possible representation-theoretic setting; the supplied record does not establish that this is the intended setting.

The requested investigation is to determine the range of cases in which Jantzen's sum formula determines the characters of simple representations, extending whatever cases are meant by the groups where determination is already complete. A precise quantified statement cannot be extracted from the supplied record: it does not specify the class of groups, the field or its characteristic, the category of representations, the version of Jantzen's sum formula, the family of simple representations under consideration, the established baseline cases, or which additional information may be used together with the formula. Accordingly, no particular group, characteristic bound, weight range, or uniqueness assertion is imposed in this draft.

The supplied dated assessment refers to low-rank cases, an \(\mathrm{SL}_5\) parameter, periodic-character methods, and hand calculations, but gives neither the parameter nor the relevant hypotheses. Those references do not define additional mathematical subproblems here.
\subsection{Short English statement}
Determine the extent to which Jantzen's sum formula can determine characters of simple representations beyond the cases already understood. The supplied record leaves the representation-theoretic setting, known cases, and meaning of determination unspecified, so this remains a qualified research question requiring clarification before a precise theorem-shaped statement can be given.
\subsection{Sources}
\begin{itemize}
\item Local registry: \texttt{data/unsolvedmath/problems.json}, with \texttt{catalog/problems.json} as the additional-record registry; supplied permanent UnsolvedMath ID \(30004016\), source problem code \texttt{catalogue entry}, title \emph{Character Determination from Jantzen's Sum Formula}. No proposer is named. The supplied proposed year is 2018.
\item Supplied report citation: \emph{Enveloping Algebras and Geometric Representation Theory} (2019), Oberwolfach Reports, pp. 3081--3146, DOI \texttt{10.4171/owr/2018/52}. \url{https://doi.org/10.4171/owr/2018/52}. The background describes Oberwolfach Reports as documenting mathematical workshops and their open-problem discussions. No problem-specific page or speaker attribution is supplied.
\item Supplied literature source: \emph{Oberwolfach Report 14/2019}. \url{https://ems.press/content/serial-article-files/46794}.
\item Supplied literature source: \emph{Coxeter Combinatorics for Sum Formulas}. \url{https://arxiv.org/abs/2107.03310}.
\item Supplied literature source: \emph{A Combinatorial Translation Principle for the Symplectic Group}. \url{https://link.springer.com/article/10.1007/s00031-021-09685-6}.
\item Literature assessment dated \texttt{2026-08-21}: this is supplied dated curation, not independently verified. The supplied status is \texttt{partially\_solved}. The assessment says that Jantzen's formula reduces the cited low-rank cases to a small number of ambiguities; later periodic-character methods and hand calculations resolve at least the \(\mathrm{SL}_5\) parameter, while no general characterization follows and higher-rank cases remain difficult. The record explicitly requests independent verification of time-sensitive status and cited claims.
\end{itemize}
\end{document}
